\subsection{An order-four element in a multiple holomorph quotient: Problem 20.108}
\label{subsec:holomorph}
\problemstatement{20.108}

For a group $G$, write $\Hol(G)$ for the normalizer of its
left regular action in the permutation group on its underlying
set, and $\NHol(G)$ for the normalizer of $\Hol(G)$.
Tsang's \kn{20.108}(a,b) asks whether a centerless group,
and specifically a finite one, can have
$T(G)=\NHol(G)/\Hol(G)$ which is not elementary abelian of
exponent two.

\begin{theorem}
\label{thm:holomorph}
For $G=(C_{11}\times C_{11})\rtimes C_5$, where a generator
of $C_5$ acts by $\operatorname{diag}(3,9)$ over $\F_{11}$,
the group $G$ is centerless and $T(G)$ contains an element
of order four.
\end{theorem}

\begin{proof}
Put $\lambda=3\in\F_{11}$, of order five, and represent
$G$ by triples with multiplication
\[
(x,y,k)(x',y',k')
 =\bigl(x+\lambda^kx',\,y+\lambda^{2k}y',\,k+k'\bigr),
\]
where $k,k'\in\Z/5\Z$. Commuting with all translations
forces a central element to have $k=0$; commuting with
$(0,0,1)$ then forces $x=y=0$. Thus $Z(G)=1$.

The unique Sylow 11-subgroup $V=\F_{11}^2$ is
characteristic. An automorphism induces a matrix $M$
on $V$ and multiplication by
$u\in(\Z/5\Z)^\times$ on $G/V$. Compatibility requires
\[
M\operatorname{diag}(\lambda,\lambda^2)
 =\operatorname{diag}(\lambda^u,\lambda^{2u})M.
\]
Equality of the eigenvalue sets gives
$\{u,2u\}=\{1,2\}$ modulo five, so $u=1$ and $M$ is
diagonal. The image of a quotient generator can have
any translation part. Conversely, the resulting maps
\[
(x,y,k)\longmapsto
\bigl(ax+b(1-\lambda^k),\,
dy+e(1-\lambda^{2k}),\,k\bigr),
\qquad a,d\ne0,
\]
are automorphisms.

A normalizer of the left regular group, after a
translation making it fix the identity, is an
automorphism. Consequently the full holomorph consists
exactly of the permutations
\[
h(x,y,k)=
\bigl(Ax+B+C\lambda^k,\,
Dy+E+F\lambda^{2k},\,k+j\bigr),
\tag{$*$}
\]
with $A,D\ne0$ and all remaining parameters arbitrary.
Define
\[
\theta(x,y,k)=(y,\lambda^{-k}x,2k),\qquad
\theta^{-1}(X,Y,K)=(\lambda^{3K}Y,X,3K).
\]
Direct composition gives
\begin{align*}
(\theta h\theta^{-1})(X,Y,K)
 =\bigl(&DX+E+F\lambda^K,\\
        &\lambda^{-j}AY+\lambda^{-j}C+
                    \lambda^{-j}B\lambda^{2K},\ K+2j\bigr).
\end{align*}
This again has form $(*)$. Since the holomorph is finite,
the resulting containment is equality, so $\theta$ normalizes it.
Moreover,
\[
\theta^2(x,y,k)=(\lambda^{-k}x,\lambda^{-2k}y,-k),
\qquad \theta^4=1.
\]
Every holomorph element acts on the final coordinate
by a translation, whereas $\theta^2$ acts by negation.
These are different functions on $\Z/5\Z$. Therefore
$\theta^2\notin\Hol(G)$ and the coset of $\theta$ has
order exactly four.
\end{proof}

Tsang's paper~\cite{tsang} singled out this group among
exceptions to its computational coverage. The proof
above uses that lead but not the paper's general
classification argument. The experimental checks
compare the explicit holomorph with GAP's automorphism
group and evaluate the normalizer identity separately.
Determining all of $T(G)$ is unnecessary for the stated
conclusion.

